\subsection{Вязко-упругий материал Кельвина–Фойгта}

\inputfig{figures/fig-kelvin-voigt}

Пружина (элемент 1) и демпфер (элемент 2) соединены параллельно:
\[
  \Sdev = \Sdev_1 + \Sdev_2 ,
  \qquad
  \dev = \dev_1 = \dev_2 .
\]
Элементы по отдельности:
\[
  \Sdev_1 = 2\mu\,\dev_1 = 2\mu\,\dev ,
  \qquad
  \Sdev_2 = 2\eta\,\dt{\dev}_2 = 2\eta\,\dt{\dev}
\]
\[
  \Longrightarrow\quad
  \boxed{\ \Sdev = 2\mu\,\dev + 2\eta\,\dt{\dev}\ }
\]
— определяющее соотношение для Кельвина–Фойгта.

\subsubsection*{Статическое нагружение}

Пусть
\[
  \Sdev = \Sdev_0\, H(t), \qquad \Sdev_0 = \mathrm{const} ,
  \qquad
  H(t) = \begin{cases} 0, & t < 0, \\ 1, & t \ge 0 \end{cases}
\]
(функция Хевисайда — мгновенно приложенная и дальше постоянная нагрузка). Тогда
\[
  \begin{cases}
    t > 0: & 2\mu\,\dev + 2\eta\,\dt{\dev} = \Sdev_0 , \\[0.3em]
    t \le 0: & \dev = 0, \quad \Sdev = 0 .
  \end{cases}
\]
Структуру девиатора $\Sdev_0$ мы не трогаем — ищем $\dev$ в виде
\[
  \dev = \Sdev_0 \,\frac{1}{2\mu}\, \gamma(t) ,
\]
где $\gamma(t)$ — скалярная функция.

\begin{warn}
Фраза перед этой подстановкой в тетради неразборчива («мы не \dots{} структуру
$\Sdev$»). Смысл: тензорная структура решения повторяет структуру нагрузки,
неизвестна только скалярная зависимость от времени.
\end{warn}

Подставляем:
\[
  \Sdev_0\gamma + \underbrace{\frac{2\eta}{2\mu}}_{\tau}\dt\gamma\,\Sdev_0 = \Sdev_0
  \quad\Longrightarrow\quad
  \tau\dt\gamma + \gamma = 1 ,
  \qquad
  \tau = \frac{\eta}{\mu} ,
\]
где $\tau$ — постоянное \textbf{время запаздывания ползучести}. Итого задача
\[
  \begin{cases}
    t > 0: & \tau\dt\gamma + \gamma = 1 , \\
    t = 0: & \gamma = 0 .
  \end{cases}
\]
Частное решение $\gamma_{\text{част}} = 1$, общее решение однородного уравнения
$\gamma_{\text{одн}} = C e^{-t/\tau}$, значит $\gamma(t) = 1 + Ce^{-t/\tau}$, и из
начального условия
\[
  \gamma(t) = 1 - e^{-t/\tau} .
\]
Тогда
\[
  \dev = \underbrace{\frac{1}{2\mu}\left(1 - e^{-t/\tau}\right)}_{\Pi(t)}\Sdev_0
  = \Pi(t)\,\Sdev_0 ,
  \qquad
  \boxed{\ \Pi(t) = \frac{1}{2\mu}\left(1 - e^{-t/\tau}\right)\ }
\]
— \textbf{функция ползучести}.

\inputfig{figures/fig-creep-function}

Свойства:
\[
  \dt\Pi(t=0) = \frac{1}{2\mu\tau} = \frac{1}{2\eta} ,
  \qquad
  \Pi(\tau) = \frac{1}{2\mu}\left(1 - e^{-1}\right) \approx \frac{2}{3}\cdot\frac{1}{2\mu} .
\]
Поэтому $\tau$ — время запаздывания ползучести: это время, за которое набирается
$\approx \tfrac23$ предельной деформации.

\subsubsection*{Произвольное нагружение: интеграл Стилтьеса}

\[
  \boxed{\ \dev = \int\limits_0^t \Pi(t-\xi)\,\dt{\Sdev}(\xi)\,d\xi
         = \int\limits_0^t \Pi(t-\xi)\,d\Sdev(\xi)\ }
\]
— интеграл Стилтьеса.

\begin{idea}
В тетради переменная интегрирования обозначена той же буквой $\tau$, что и время
запаздывания. Здесь она переименована в $\xi$, чтобы не путать: $\tau$ всюду
дальше — характерное время материала.

Форма Стилтьеса удобна именно потому, что допускает скачки $\Sdev$: в точке
скачка $d\Sdev$ даёт конечный вклад, тогда как $\dt{\Sdev}$ там — дельта-функция.
\end{idea}

Проверим на примерах.

\medskip
\textbf{1)} $\Sdev = \Sdev_0 H(t)$. Тогда $d\Sdev = \Sdev_0\,\delta(\xi)d\xi$ и
\[
  \dev(t) = \Pi(t)\,\Sdev_0 ,
\]
то есть результат согласуется с прямым решением.

\medskip
\textbf{2)} Импульс нагрузки
\[
  \Sdev = \Sdev_0 f(t) = \Sdev_0 H(t) - \Sdev_0 H(t - t^*) :
\]
\[
  \dev(t) = \Sdev_0\bigl(\Pi(t) - \Pi(t - t^*)\bigr) = \Sdev_0\,\Lambda(t) .
\]
При этом $\Pi(t)$ всегда больше, чем $\Pi(t-t^*)$, — функция ползучести растёт,
поэтому $\Lambda(t) \ge 0$: после снятия нагрузки деформация не исчезает мгновенно,
а убывает по экспоненте (обратная ползучесть).

\inputfig{figures/fig-creep-pulse}

\subsection{Обобщённый материал Кельвина–Фойгта}

Последовательное соединение $N$ штук моделей Кельвина–Фойгта.

\inputfig{figures/fig-kv-generalized}

По аксиомам последовательного соединения
\[
  \dev = \sum_{k=1}^{N} \dev_k ,
  \qquad
  \Sdev = \Sdev_1 = \Sdev_2 = \dots = \Sdev_N ,
\]
а каждое звено даёт
\[
  \dev_k = \int\limits_0^t \frac{1}{2\mu_k}\left(1 - e^{-(t-\xi)/\tau_k}\right) d\Sdev(\xi) ,
  \qquad
  \tau_k = \frac{\eta_k}{\mu_k} .
\]
Складывая:
\[
  \dev = \int\limits_0^t \underbrace{\left[\sum_{k=1}^{N}\frac{1}{2\mu_k}
         \left(1 - e^{-(t-\xi)/\tau_k}\right)\right]}_{\Pi(t-\xi)} d\Sdev(\xi)
       = \int\limits_0^t \Pi(t-\xi)\,d\Sdev(\xi) .
\]

Предельные значения:
\[
  t = 0 \ \Longrightarrow\ \dt\Pi(0) = \sum_{k=1}^{N}\frac{1}{2\eta_k} ,
  \qquad
  t \to \infty \ \Longrightarrow\ \Pi(t) \to \sum_{k=1}^{N}\frac{1}{2\mu_k} .
\]

\begin{warn}
Следующая фраза в тетради читается плохо: «в н. момент напряжения только
в \dots, в \dots{} напряжения только в \dots». По физике модели: сразу после
скачка нагрузки деформация равна нулю, скорость деформации конечна только за счёт
демпферов — нагрузку несут \emph{демпферы}; при $t\to\infty$ скорости обнуляются
и всю нагрузку несут \emph{пружины}.
\end{warn}

При $N = 2$, $\mu_1 = \mu_2$, $\tau_2 \gg \tau_1$ график $\Pi(t)$ имеет излом
около $t = \tau_1$ (на уровне $1/2\mu_1$) и выходит на $1/\mu_1$ при $t\sim\tau_2$;
реальный график — сглаженный.

\inputfig{figures/fig-kv-two-times}
